% Harmonics against level, load and bias (triode-distortion): the set-up a person would wire, semantically the deck's.
%   VBBi bbi 0 <ebb>, RPi bbi ami <rp>, VAMi ami ai 0   the supply, the plate resistor, the anode ammeter   -> source E_bb (set), resistor R_p, ammeter I_a (measured)
%   VKi ki 0 <vk>                                       the cathode on an ideal bypass: Rk times the mean cathode current, iterated -> source (set)
%   VINi ini 0 SIN(0 <drive> 1k), RG1i ini 0 <rg1>, VGMi ini gi 0   the drive, the grid resistor, the grid ammeter -> source v_in (set)
%   CCi ami oi 10u, RSi oi 0 <rs>                       the coupling capacitor and the next grid's resistor  -> voltmeter v_out (measured)
%   XTi ai gi ki DUT; .tran 30 periods, the last ten read; a quiescent .op with the real Rk starts the iteration
%   Fifteen copies, one per drive level (half-octave steps from the drive that gives 2 V rms out), drawn as one.
\documentclass[tikz, border=10pt]{standalone}
\usepackage{bench}
\begin{document}
\begin{tikzpicture}
\benchtitle{0,8.6}{Harmonics against level, load and bias}
\def\gnd{1.0}
% the triode
\node[triode, scale=1.35] (T) at (4,3.8) {};
\node[anchor=west, font=\small, align=left] at ($(T.anode)!0.62!(T.cathode)+(0.75,0)$) {DUT};
% copy : the anode resistor up to the rail, the anode ammeter under it
\coordinate (am) at ($(T.anode)+(0,1.1)$);
\ammeter{T.anode}{am}{$I_a$}
\draw (am) to[R, l_={$R_p$}] ++(0,1.6) coordinate (top);
\coordinate (vtop) at ($(top)+(6.6,0)$);
\draw (top) -- (vtop);
\draw[set] (vtop) to[V, l={$E_{bb}$}] (vtop |- 0,\gnd) node[ground, black]{};
\nodename{anchor=south}{$(top)+(1.0,0.05)$}{bb}
\nodename{anchor=east}{$(am)+(-0.15,0.15)$}{a}
% the cathode resistor, and its bypass
\draw (T.cathode) to[R, l={$R_k$}] (T.cathode |- 0,\gnd) node[ground]{};
\draw (T.cathode) -- ++(1.2,0) coordinate (ck) to[C, l={$C_k$}] (ck |- 0,\gnd) node[ground]{};
\nodename{anchor=east}{$(T.cathode)+(-0.1,-0.25)$}{k}
% the grid: its resistor to ground and the source
\coordinate (g) at (T.control);
\draw (g) -- ++(-0.6,0) coordinate (gj);
\draw (gj) to[R, l_={$R_{g1}$}] (gj |- 0,\gnd) node[ground]{};
\draw (gj) -- ++(-0.3,0) to[R, l_={$R_{src}$}] ++(-1.8,0) coordinate (vin);
\draw[set] (vin) to[V, l_={$v_{in}$}] (vin |- 0,\gnd) node[ground, black]{};
\nodename{anchor=south}{$(g)+(-0.2,0.05)$}{g}
\nodename{anchor=south east}{$(vin)+(-0.05,0.05)$}{in}
% the output: through Cc into the next grid's resistor
\coordinate (oj) at ($(am)+(0,0.0)$);
\draw (am) -- ++(1.6,0) to[C, l={$C_c$}] ++(1.4,0) coordinate (o);
\draw (o) to[R, l_={$R_s$}] (o |- 0,\gnd) node[ground]{};
\nodename{anchor=south}{$(o)+(0.3,0.05)$}{o}
% the voltmeter on o
\coordinate (vo) at ($(o)+(1.2,0)$);
\draw (o) -- (vo);
\voltmeter{vo |- 0,\gnd}{vo}{$v_{out}$}
\draw (vo |- 0,\gnd) node[ground]{};
% the definitions
\definitions{0,-0.3}{%
  the chart's middle cell, and a ladder around it: $v_{out} = 0.1$, 0.3, 1, 2, 3, 5, 10, 20, 30, 50 V rms;\quad $R_p$ and $R_k$ each $\times\,0.5$, 1, 2\\[2pt]
  $C_k$: bypassed (ideal), 25 \textmu F, and none (local feedback);\quad $R_{src} = 0$, 10 k, 100 k\ohm\\[2pt]
  read $v_{out}$ over the last ten periods: THD, H2 to H5 by Goertzel, the sign of H2 (the phase of the second harmonic against the fundamental)\\[2pt]
  model against model: the sheets print distortion only for output stages and the RC chart's 5\,\% point%ENDDEFS
}
\end{tikzpicture}
\end{document}
